%% notes-tex/024-perturb-seq-pilot/sections/5-checkpoints.tex
%% SOURCE: hand-authored synthesis. Sequencing prices and cell counts are from
%%         experiments/024-perturb-seq-costing/scripts/pilot_options.py and
%%         uiuc_core_data.py; the counting and core-service facts are Sec. 5.5 of
%%         the costing review; the three checks are Sec. 3 of this document.

\section{Running It\secstatus{todo}}
\label{sec:checkpoints}

\subsection{What Each Run Must Return\secstatus{todo}}
\label{sec:returns}

A run is finished when it has produced these, and not when it has produced data.

\notesec{\org{S.\ cerevisiae}, one 5$'$ channel on the existing CRISPRi library}
\begin{enumerate}[leftmargin=1.6em,itemsep=2pt,topsep=2pt]
  \item $q$, the fraction of quality-controlled cells carrying exactly one
    identified guide. This is the number the run exists for.
  \item The targeted gene's own transcript in the cells carrying its guide,
    against cells carrying other guides. Interference should reduce it, and this
    is the positive control whose absence went unnoticed in \cref{sec:nadal}.
  \item Messenger RNA molecules and genes per cell, and the fraction of reads
    that are ribosomal.
  \item Guide representation: how many of the library's guides appear at all, and
    how unevenly.
\end{enumerate}

\notesec{\org{P.\ kudriavzevii}, one channel, several media, no library}
\begin{enumerate}[leftmargin=1.6em,itemsep=2pt,topsep=2pt]
  \item Whether the wall yields, read as cells recovered against cells loaded.
  \item Messenger RNA molecules and genes per cell.
  \item Ribosomal fraction, which is the largest single lever on the cost of
    every later experiment and is a proportion, so it is answerable at any
    sequencing depth.
  \item Whether the media separate, which is both a biological result and the
    evidence that condition labels survive the protocol.
\end{enumerate}

\subsection{Three Things to Get Right at the Bench\secstatus{todo}}
\label{sec:counting}

\notesec{Count the cells on an instrument that can see yeast}
\org{S.\ cerevisiae} runs about 3.5 to 5\,\textmu m
across\external{vendor sizing notes; no primary in the mirror}, at or below the
stated floor of most brightfield counters, and 10x's own guidance is that cells
under 10\,\textmu m ``may be difficult to distinguish from debris'' and that a
fluorescent viability stain rather than trypan blue is ``strongly
recommended''\external{10x CG00053}. Two instruments are validated on this
organism by their makers, the Logos LUNA-FX7 and the ChemoMetec
NC-203\external{Logos and ChemoMetec application notes}. The Bio-Rad TC20
specifies a 6\,\textmu m floor\external{Bio-Rad TC20 brochure}, which excludes
yeast. The count sets the droplet loading rate and therefore the doublet rate,
and the core named it as the single most critical step.

\notesec{Replicate one genotype across two independent wells} Not a second
channel and not a technical resequencing. Two wells, or two cartridges, with
wild-type cells in each, so that a profile can be compared against itself with a
reference from its own batch. This is the check that would have caught the
published screen, it costs a handful of wells, and without it a null result
later cannot be told from a broken assay.

\notesec{Sequence twice, shallow then deep} A MiSeq run confirms the chemistry
worked, and the questions it answers are proportions, so they need almost no
depth. A NovaSeq X 10B split lane then buys 600 million read pairs for \$1{,}160,
which is less than the MiSeq run costs and 24 times the reads. There is no
earlier checkpoint than sequencing something: complementary DNA traces and
library concentration do not predict whether barcoding
worked~\citep{brettnerUltraHighthroughputMassively2024}.

\subsection{What Happens Next, Under Each Outcome\secstatus{todo}}
\label{sec:next}

\notesec{If $q$ is high in \org{S.\ cerevisiae}} The library needs no rebuilding.
Run a focused single-guide pilot on a metabolic sub-library at 100 cells per
gene, and add growth media to it through the plate route, where they are close
to free.

\notesec{If $q$ is near zero} The 5$'$ scaffold primer does not engage a
polymerase III transcript. Build the polyadenylated guide barcode and validate
it on about ten arrayed strains with known knockdown phenotypes against bulk RNA
sequencing of the same strains. Ground truth is only available while the design
is arrayed, since a pool cannot say what was in a given cell.

\notesec{If the \org{P.\ kudriavzevii} wall does not yield} The failure is in
sample preparation rather than in the platform, and the next attempt varies the
digestion rather than the sequencing. Nothing else in the plan depends on it,
which is why it is worth running now.

\notesec{If both work} The two tracks converge on one question, which is whether
to spend the next budget on more perturbations per cell or on more growth media.
The arithmetic favors media, and the model that would consume the data needs
both, so the first crossed design is a small metabolic panel across a handful of
conditions rather than a genome-scale screen in one.

\begin{keybox}
\textbf{Two runs, two people, about \$8{,}300, and every later decision depends
on what they return.} The \org{S.\ cerevisiae} channel decides whether a
construction project stands between us and a screen. The
\org{P.\ kudriavzevii} channel decides whether the assay works on that organism
at all. Build the three quality checks of \cref{sec:nadal} into both, because
the published screen they come from is the design this plan would otherwise have
copied.
\end{keybox}
